\documentclass[12pt,a4paper]{article}

% ===============================
% PAQUETES BÁSICOS
% ===============================
\usepackage[spanish, es-noshorthands]{babel}   % Idioma español
\usepackage[utf8]{inputenc}                   % Codificación UTF-8
\usepackage[T1]{fontenc}                      % Caracteres con acentos correctos
\usepackage{lmodern}                          % Fuente moderna
\usepackage[a4paper, margin=2.5cm]{geometry}  % Márgenes reducidos
\usepackage{setspace}
\onehalfspacing                               % Interlineado 1.5

% ===============================
% IMÁGENES Y GRÁFICOS
% ===============================
\usepackage{graphicx}     % Insertar imágenes
\usepackage{float}        % Control de posición (H)
\usepackage{caption}
\usepackage{subcaption}

% ===============================
% TABLAS
% ===============================
\usepackage{booktabs}     % Líneas elegantes
\usepackage{multirow}
\usepackage{tabularx}
\usepackage{array}

% ===============================
% MATEMÁTICAS
% ===============================
\usepackage{amsmath, amssymb, amsfonts}
\usepackage{siunitx}

% ===============================
% OTROS ÚTILES
% ===============================
\usepackage{xcolor}
\usepackage[hidelinks]{hyperref}
\usepackage{enumitem}

% ===============================
% INFORMACIÓN DEL DOCUMENTO
% ===============================
\title{Título del Artículo en Español}
\author{Nombre del Autor}
\date{\today}

% ===============================
% DOCUMENTO
% ===============================
\begin{document}
	
	\maketitle
	
	\begin{abstract}
		Este es el resumen del artículo. Explica brevemente el objetivo, el contexto y los resultados principales.
	\end{abstract}
	
	\section{Introducción}
	
	Motivación
	
	¿Para que usar topic modeling en humanidades?
	Tengo 500 articulos. Como explorar sus temáticas. 
	Existen algunas herramientas en humanidades digitales. Una de estas es tm.
	¿Que es?
	Dos tipos
	LDA y Bertopic
	Ejemplo resultado LDA
	Ejemplo resultado BERTopic
	Aplicación práctica en python
	
	
	
	
	Ejemplo de texto en español.  
	Ejemplo de cita en formato IEEE: como se indica en \cite{autor2020}, el análisis de datos es esencial en la investigación moderna.  
	También puede citarse varias fuentes a la vez \cite{gomez2019,lopez2021}.
	
	\section{Ecuaciones}
	
	Ejemplo de ecuación:
	\begin{equation}
		E = mc^2
	\end{equation}
	
	\section{Tablas}
	
	\begin{table}[H]
		\centering
		\caption{Ejemplo de tabla con booktabs}
		\begin{tabular}{lcc}
			\toprule
			Variable & Media & Desviación estándar \\
			\midrule
			X & 10.2 & 1.3 \\
			Y & 8.7 & 0.9 \\
			\bottomrule
		\end{tabular}
	\end{table}
	
	\section{Imágenes}
	
	\begin{figure}[H]
		\centering
		%\includegraphics[width=0.6\textwidth]{ejemplo-imagen.png}
		\caption{Ejemplo de imagen insertada.}
	\end{figure}
	
	\section{Conclusiones}
	
	Aquí van las conclusiones del trabajo.  
	
	% ===============================
	% BIBLIOGRAFÍA - FORMATO IEEE
	% ===============================
	\bibliographystyle{IEEEtran}     % Estilo de citas IEEE
	\bibliography{bibliografia}      % Nombre del archivo .bib (sin extensión)
	
\end{document}
